\documentclass[11pt,a4paper]{article}
\usepackage[T1]{fontenc}
\usepackage[utf8]{inputenc}
\usepackage{amsmath,amssymb,amsthm}
\usepackage[margin=1in]{geometry}
\usepackage[colorlinks=true,linkcolor=blue,urlcolor=blue]{hyperref}
\theoremstyle{plain}
\newtheorem{theorem}{Theorem}
\newtheorem{lemma}[theorem]{Lemma}
\newtheorem{proposition}[theorem]{Proposition}
\newtheorem{corollary}[theorem]{Corollary}
\theoremstyle{definition}
\newtheorem{remark}[theorem]{Remark}

\title{The empirical recurrence for OEIS A209853, proved by transfer matrix}
\author{Adrian Perez Fontelles\\ \small Independent researcher}
\date{31 August 2026}
\begin{document}
\maketitle

\begin{abstract}
OEIS A209853 counts $(n+1)\times4$ arrays over $\{0,\dots,3\}$ subject to a condition imposed on
every $2\times2$ subblock, and carries an empirical recurrence of order $68$ contributed
by R.~H.~Hardin. It is true, and it is not an empirical matter at all. The condition is local
to each subblock, so the rows of an admissible array are the vertices of a finite digraph
on $S=256$ vertices whose edges are the admissible consecutive pairs, and the count is a
number of walks: $a(n)=\tfrac1{4}\mathbf{1}^{\!\top}M^{n}\mathbf{1}$. Writing $q$
for the characteristic polynomial of the conjectured recurrence, the recurrence holds for all
$n$ exactly when $\mathbf{1}^{\!\top}M^{j}q(M)\mathbf{1}=0$ for every $j\ge0$, and the
Cayley--Hamilton theorem bounds that to $j<S$. It is therefore a finite computation in exact
integer arithmetic, and it comes out zero.
\end{abstract}

\noindent\small 2020 Mathematics Subject Classification. 05A15, 05B45, 15B36.\normalsize

\section{The sequence and the conjecture}

OEIS A209853 is ``1/4 the number of (n+1)X4 0..3 arrays with every 2X2 subblock having one or four distinct clockwise edge differences.''. It has offset $1$ and begins
\[
1308, 43091, 1386131, 45071990, 1460712837, 47417927307, 1538049629157, 49909370504583,\ \dots
\]
The entry carries the following comment:
\begin{quote}\small
Empirical: a(n) = 20*a(n-1) +963*a(n-2) -12793*a(n-3) -312156*a(n-4) +3571443*a(n-5) +49605457*a(n-6) -553758368*a(n-7) -4457972689*a(n-8) +52612881372*a(n-9) +243417930360*a(n-10) -3275949406092*a(n-11) -8240107440663*a(n-12) +140666492498369*a(n-13) +160116943902217*a(n-14) -4326535316148209*a(n-15) -817021352845345*a(n-16) +98006065095893985*a(n-17) -49203318829660909*a(n-18) -1668148030882549725*a(n-19) +1723999493198171149*a(n-20) +21637008392085651838*a(n-21) -31773088576337392703*a(n-22) -215843167651091106385*a(n-23) +399719679271900244237*a(n-24) +1664018938648323383901*a(n-25) -3697329744805424755681*a(n-26) -9917310558011937270066*a(n-27) +25982170466786070131093*a(n-28) +45447978346529331965091*a(n-29) -141117691880001505315599*a(n-30) -157856587989686977889953*a(n-31) +598069670716001721149653*a(n-32) +401968717089020637279250*a(n-33) -1987776050324662023234413*a(n-34) -686660571916972008322547*a(n-35) +5191091867684820767344742*a(n-36) +522733262632635990820808*a(n-37) -10646069432263823711177799*a(n-38) +915809068137700417207761*a(n-39) +17098571599529269338686083*a(n-40) -3943750897982196417975435*a(n-41) -21398685175106614262349856*a(n-42) +7368362612864864601239602*a(n-43) +20709879068969481472447106*a(n-44) -9015916934517168198203128*a(n-45) -15335247867307455505750474*a(n-46) +7824208837288323351242620*a(n-47) +8561662593050307213736807*a(n-48) -4920492104705864543228925*a(n-49) -3532336472390984912564596*a(n-50) +2242642650064493959464073*a(n-51) +1047216070797663839803065*a(n-52) -731775898101400159982255*a(n-53) -214107873494035007863503*a(n-54) +167132111313457973560540*a(n-55) +28238564794974400827975*a(n-56) -25825554003258952415666*a(n-57) -2098980376078516736024*a(n-58) +2571400424053876584470*a(n-59) +52829982196209769868*a(n-60) -153859888970921746848*a(n-61) +2954600261893327248*a(n-62) +4977837135363185688*a(n-63) -211003278535383648*a(n-64) -72160266595102240*a(n-65) +3672896320734208*a(n-66) +335148663461120*a(n-67) -16066999343104*a(n-68)
\end{quote}
As of the ``Last modified'' line on the live entry (Oct 03 2025 15:33:11, revision 9) the statement is
still recorded as empirical, and nothing on the entry records it as settled either way.

Written out, the claim is
\begin{equation}\label{eq:conj}
a(n)\;=\;20\,a(n-1) + 963\,a(n-2) - 12793\,a(n-3) - 312156\,a(n-4) + 3571443\,a(n-5) + 49605457\,a(n-6) - 553758368\,a(n-7) - 4457972689\,a(n-8) + 52612881372\,a(n-9) + 243417930360\,a(n-10) - 3275949406092\,a(n-11) - 8240107440663\,a(n-12) + 140666492498369\,a(n-13) + 160116943902217\,a(n-14) - 4326535316148209\,a(n-15) - 817021352845345\,a(n-16) + 98006065095893985\,a(n-17) - 49203318829660909\,a(n-18) - 1668148030882549725\,a(n-19) + 1723999493198171149\,a(n-20) + 21637008392085651838\,a(n-21) - 31773088576337392703\,a(n-22) - 215843167651091106385\,a(n-23) + 399719679271900244237\,a(n-24) + 1664018938648323383901\,a(n-25) - 3697329744805424755681\,a(n-26) - 9917310558011937270066\,a(n-27) + 25982170466786070131093\,a(n-28) + 45447978346529331965091\,a(n-29) - 141117691880001505315599\,a(n-30) - 157856587989686977889953\,a(n-31) + 598069670716001721149653\,a(n-32) + 401968717089020637279250\,a(n-33) - 1987776050324662023234413\,a(n-34) - 686660571916972008322547\,a(n-35) + 5191091867684820767344742\,a(n-36) + 522733262632635990820808\,a(n-37) - 10646069432263823711177799\,a(n-38) + 915809068137700417207761\,a(n-39) + 17098571599529269338686083\,a(n-40) - 3943750897982196417975435\,a(n-41) - 21398685175106614262349856\,a(n-42) + 7368362612864864601239602\,a(n-43) + 20709879068969481472447106\,a(n-44) - 9015916934517168198203128\,a(n-45) - 15335247867307455505750474\,a(n-46) + 7824208837288323351242620\,a(n-47) + 8561662593050307213736807\,a(n-48) - 4920492104705864543228925\,a(n-49) - 3532336472390984912564596\,a(n-50) + 2242642650064493959464073\,a(n-51) + 1047216070797663839803065\,a(n-52) - 731775898101400159982255\,a(n-53) - 214107873494035007863503\,a(n-54) + 167132111313457973560540\,a(n-55) + 28238564794974400827975\,a(n-56) - 25825554003258952415666\,a(n-57) - 2098980376078516736024\,a(n-58) + 2571400424053876584470\,a(n-59) + 52829982196209769868\,a(n-60) - 153859888970921746848\,a(n-61) + 2954600261893327248\,a(n-62) + 4977837135363185688\,a(n-63) - 211003278535383648\,a(n-64) - 72160266595102240\,a(n-65) + 3672896320734208\,a(n-66) + 335148663461120\,a(n-67) - 16066999343104\,a(n-68)
\end{equation}
for all $n>68$.

\section{The condition is local, so the rows form a digraph}

Write an array of the shape in question as its list of rows
$r^{(1)},r^{(2)},\dots$, each an element of
$V=\{0,1,\dots,3\}^{4}$, so $|V|=S=256$. A $2\times2$ subblock of the array
occupies two consecutive rows and two consecutive columns; if the two rows are $r$
and $s$ (in that order) and the two columns are numbered $j,j+1$,
the subblock is
\[
\begin{pmatrix}\alpha&\beta\\\gamma&\delta\end{pmatrix}=\begin{pmatrix}r_j&r_{j+1}\\s_j&s_{j+1}\end{pmatrix}.
\]
The entry's requirement on that subblock is
\begin{equation}\label{eq:loc}
\#\{\beta-\alpha,\ \delta-\beta,\ \gamma-\delta,\ \alpha-\gamma\}\in\{1,4\},
\end{equation}
a condition on the four entries alone.

For $r,s\in V$ declare $r\to s$ an edge of a digraph $G$ when, for every
$1\le j\le 4-1$,
\begin{itemize}
\item the four entries above satisfy \eqref{eq:loc};

\end{itemize}
Let $M\in\{0,1\}^{S\times S}$ be the adjacency matrix of $G$ and $\mathbf{1}$ the
all-ones vector.

\begin{lemma}\label{lem:walk}
The arrays counted by the entry with $L$ rows are exactly the walks
$r^{(1)}\to r^{(2)}\to\cdots\to r^{(L)}$ of length $L-1$ in $G$. Consequently
\[
a(n)\;=\;\tfrac1{4}\mathbf{1}^{\!\top}M^{n}\mathbf{1} .
\]
\end{lemma}

\begin{proof}
An array with $L$ rows and $4$ columns is precisely a sequence
$r^{(1)},\dots,r^{(L)}$ of elements of $V$; conversely any such sequence is an array.
Every $2\times2$ subblock lies in two consecutive rows $r^{(t)},r^{(t+1)}$ and two
consecutive columns, so the array satisfies the entry's condition on all of its subblocks
if and only if each consecutive pair $\bigl(r^{(t)},r^{(t+1)}\bigr)$ satisfies it for
every column pair --- which is exactly the edge relation of $G$. The number of walks of
length $L-1$, summed over all starting and ending vertices, is
$\mathbf{1}^{\!\top}M^{L-1}\mathbf{1}$, since $(M^{k})_{r,s}$ counts the walks of
length $k$ from $r$ to $s$. The entry's index $n$ corresponds to $L=n+1$ rows. The entry counts $\tfrac1{4}$ of those arrays, a constant factor that passes through every step below unchanged.
\end{proof}

\section{The criterion}

Let $q(t)=t^{68}-\sum_i c_i\,t^{68-i}$ be the characteristic polynomial of
\eqref{eq:conj}, with the $c_i$ read off that equation.

\begin{lemma}\label{lem:crit}
For every $n$ with $n+1-1\ge68$,
\[
a(n)-\sum_i c_i\,a(n-i)\;=\;\tfrac1{4}\mathbf{1}^{\!\top}M^{\,n+1-1-68}\,q(M)\,\mathbf{1} .
\]
Hence \eqref{eq:conj} holds from that point on if and only if
$\mathbf{1}^{\!\top}M^{j}q(M)\mathbf{1}=0$ for every $j\ge0$; and it is enough to check
$j=0,1,\dots,S-1$.
\end{lemma}

\begin{proof}
By Lemma~\ref{lem:walk}, $a(n-i)=\tfrac1{4}\mathbf{1}^{\!\top}M^{\,n+1-1-i}\mathbf{1}$, so
\[
a(n)-\sum_i c_i a(n-i)
=\tfrac1{4}\mathbf{1}^{\!\top}\Bigl(M^{m}-\sum_i c_i M^{m-i}\Bigr)\mathbf{1}
=\tfrac1{4}\mathbf{1}^{\!\top}M^{\,m-68}\,q(M)\,\mathbf{1}
\]
with $m=n+1-1$, which is the stated identity; the scalar $\tfrac1{4}$ never vanishes, so it does not affect whether the left side is zero, and the equivalence follows by letting
$m-68=j$ run over $\ge0$. For the last clause put $w=q(M)\mathbf{1}$. By the
Cayley--Hamilton theorem $M^{S}$ is an integer combination of $I,M,\dots,M^{S-1}$, so
$\mathbf{1}^{\!\top}M^{j}w$ for $j\ge S$ is the corresponding combination of the earlier
values; if those all vanish, so do all the rest.
\end{proof}

\section{The computation}

\begin{theorem}
The recurrence \eqref{eq:conj} holds for every $n>68$, which is the statement of
Section~1.
\end{theorem}

\begin{proof}
The digraph was constructed from \eqref{eq:loc} by a depth-first scan across the $4$
columns: the edge condition is a conjunction of one constraint per consecutive column
pair, so the successors of a row $r$ are enumerated column by column without testing all
$S^2$ pairs. The vector $w=q(M)\mathbf{1}\in\mathbb{Z}^{S}$ was then formed by $68$
matrix--vector products in exact integer arithmetic, and $\mathbf{1}^{\!\top}M^{j}w$ was
evaluated for $j=0,1,\dots$, again exactly. Every one of those integers vanishes from the
index corresponding to $n=68+1$ onwards, and the run continued until $w$ itself was the
zero vector, which settles every larger $j$ at once. By Lemma~\ref{lem:crit} the recurrence
holds for every $n>68$.
\end{proof}

\begin{remark}
The argument gives more than the recurrence: it shows the sequence is $C$-finite with
characteristic polynomial dividing that of $M$, so it has a rational generating function
whose denominator divides $\det(I-xM)$. The conjectured recurrence is one consequence; the
minimal one is a divisor of $q$.
\end{remark}

\section{Verification}

Three checks, each independent of the algebra above.

First, the digraph was built directly from the entry's stated condition and the walk counts
$\tfrac1{4}\mathbf{1}^{\!\top}M^{L-1}\mathbf{1}$ compared against the entry's DATA at
every one of the $14$ published indices; the values agree exactly. That is what ties
the digraph to the sequence: a model that reproduced only some of the published terms would
be the wrong model, and would have been rejected. In particular it is what rules out a
misreading of the entry's wording, since a different reading of the condition gives a
different digraph and different counts.

Second, the recurrence \eqref{eq:conj} was evaluated on those published terms in exact
integer arithmetic, with no matrices involved, and vanishes wherever it applies.

Third, the annihilation test of Lemma~\ref{lem:crit} was run in exact integer arithmetic
until the working vector was identically zero, not to a sampled prefix, and returned zero
throughout.

\begin{thebibliography}{9}
\bibitem{oeis} The OEIS Foundation, \emph{The On-Line Encyclopedia of Integer
Sequences}, \url{https://oeis.org}, 2026. Sequence A209853.
\bibitem{stanley} R.~P.~Stanley, \emph{Enumerative Combinatorics}, Volume 1, 2nd ed.,
Cambridge University Press, 2011. (Transfer-matrix method, Section 4.7.)
\bibitem{fs} P.~Flajolet and R.~Sedgewick, \emph{Analytic Combinatorics}, Cambridge
University Press, 2009.
\bibitem{horn} R.~A.~Horn and C.~R.~Johnson, \emph{Matrix Analysis}, 2nd ed., Cambridge
University Press, 2013.
\end{thebibliography}

\end{document}
